% !TEX root = ../main.tex
\chapter{Differentiable Manifolds}\label{ch:c04}
\begin{paracol}{2}

\begin{sources}
\begin{tabularx}{\linewidth}{@{}l l L@{}}
\toprule
\textbf{Source} & \textbf{Length} & \textbf{Content}\\
\midrule
\seg{4}{1}{Lecture 4: Differentiable Manifolds} & 1:00:24 & Why a topological manifold is not enough to speak of differentiability; the strategy through chart representatives and why it fails;
compatible charts and $C^k$, $C^\infty$, $C^\omega$, complex atlases; Whitney's theorem and smooth manifolds; student questions on the trajectory of a piece of chalk and on $C^2$ curves; isomorphisms of sets, topological spaces and vector
spaces, diffeomorphisms; the number of smooth structures on a topological manifold in each dimension.\\
\seg{4}{2}{Tutorial 4: Differentiable Manifolds} & 0:35:23 & Solutions of the sheet by a tutor (not Schuller): fifteen true/false statements on differentiability and manifolds; an atlas on $\R$ that is not smooth, and how to repair it; an atlas of squares on $\R^2$;
partial derivatives as maps; chart representatives of a curve and of a function and a chain rule computation.\\
\bottomrule
\end{tabularx}

\smallskip
\textbf{Transcripts} \file{materials/transcripts/C04-V01.txt} and \file{C04-V02.txt} (YouTube captions).
\textbf{Sheet 4:} see \url{https://tales.mbivert.com/on-heraeus-winter-school-tutorials/}.
\end{sources}

\begin{reading}[Where this chapter sits in the other notes]
\emph{Geometric Anatomy of Theoretical Physics} notes: ``Differentiable structures: definition and classification'' (chapter 7: compatible charts, $C^k$ atlases,
the maximal smooth atlas, smooth maps, diffeomorphisms, and the classification of smooth structures with the results quoted at the end of this lecture); the topological manifolds it builds on are chapter 6, and the vector spaces used
in the next lecture are chapter 8 (this book's chapter~\ref{ch:c03}). Lecture 5 (tangent spaces) continues with ``Construction of the tangent bundle'' (chapter 10 there).
\end{reading}

% ==================================================================
\section{The problem: differentiability needs more structure}\label{sec:c04-problem}

Welcome to the fourth lecture \ts{4}{1}{0:00}. So far topological manifolds: the topology gives the definition of continuity of a curve, and if one thinks of a curve as the trajectory of a particle in classical physics, one wants to require that it does not make jumps; this has been achieved.
Now one wants more: to associate with the curve a \emph{velocity} at every point. If the curve has a corner, there is no unique tangent vector there (very roughly speaking); the first curve is intuitively differentiable, the second is not. This needs to be made precise
\ts{4}{1}{0:20}\ts{4}{1}{0:41}\ts{4}{1}{1:02}\ts{4}{1}{1:23}.

The question is whether the structure $(M,\mathcal O)$ of a topological manifold of dimension $d$ is enough to talk about differentiability of curves. ``The answer will be a resounding no'': there is not enough information in it, one needs to make more choices, and this lecture
will find out how much choice there is \ts{4}{1}{1:23}\ts{4}{1}{1:43}\ts{4}{1}{2:04}.
Already now: in dimension $1,2,3$ there are choices but they are all equivalent, so there is essentially only one way to talk about differentiability of curves; in dimension $5,6,7,\dots$ there is in each dimension a finite amount of choice (``you say, well, you have 48 choices: write it down,
do physics with it, make experiments and you will find out which one it is''); but unfortunately, there is a theorem (quoted at the end) that in \emph{four} dimensions, the dimension of spacetime, there is an uncountable number of possible choices. ``Do you still think it is obvious which one to choose? It is not.''
Hence it is important to figure out what is going on \ts{4}{1}{2:04}\ts{4}{1}{2:25}\ts{4}{1}{2:45}\ts{4}{1}{3:05}.

\paragraph{Goal.} To find what structure to add to a topological manifold to talk about differentiability of \emph{curves} $\gamma\colon\R\to M$, of \emph{functions} on the manifold (``that is just the terminology for a map from the manifold into the real numbers'', say a temperature distribution), and most
generally of \emph{maps from one manifold to another}; all of this will be needed. The topology is in the background and will not always be specified \ts{4}{1}{3:05}\ts{4}{1}{3:26}\ts{4}{1}{3:49}\ts{4}{1}{4:10}.

\section{The strategy and why it fails}\label{sec:c04-strategy}

The strategy was already indicated in lecture 2 (\cref{sec:c02-philosophy}). If one looked at a curve from $\R$ to $\R^d$ rather than into the manifold, one would know what differentiability means: the undergraduate definition. So choose a chart $(U,x)$ (it does not cover the whole manifold, but start with the
portion of the curve in the chart domain; since chart domains cover everything, one then composes the global notion from the local one) and look at $x\circ\gamma$, a curve from $\R$ to $\R^d$ whose differentiability one can judge, with ``your favourite definition'' (the derivative with respect to the curve parameter, componentwise)
\ts{4}{1}{4:30}\ts{4}{1}{4:50}\ts{4}{1}{5:11}\ts{4}{1}{5:31}\ts{4}{1}{5:52}\ts{4}{1}{6:14}. The idea is to \emph{lift} the undergraduate notion of differentiability of a curve on $\R^d$ to a notion of ``manifold differentiability'' of a curve on $M$ \ts{4}{1}{6:14}\ts{4}{1}{6:37}\ts{4}{1}{6:59}\ts{4}{1}{7:20}.

The question already met last time: is this \emph{well defined}? The curve does not know anything about my looking at a particular chart, and it would be a very bad idea to assign it a property that depends on my taste \ts{4}{1}{7:20}\ts{4}{1}{7:40}.
Take another chart $(V,y)$ with $U\cap V\neq\emptyset$ (if the intersection is empty the two charts cannot be incompatible: that portion of the curve has nothing to do with the other). On the overlap one chart map might say that the curve is differentiable and the other that it is not, ``a catastrophe, then our definition is no good''.
So look at the overlap $U\cap V$, mapped by $x$ and by $y$ into $\R^d$, giving the representatives $x\circ\gamma$ and $y\circ\gamma$ \ts{4}{1}{8:01}\ts{4}{1}{8:22}\ts{4}{1}{8:42}\ts{4}{1}{9:03}\ts{4}{1}{9:23}\ts{4}{1}{9:43}.
If $x\circ\gamma$ is undergraduate differentiable, is $y\circ\gamma$? One has the chart transition map $y\circ x^{-1}$ and
\[
  y\circ\gamma=(y\circ x^{-1})\circ(x\circ\gamma)
\]
(inserting the identity $x^{-1}\circ x$; the composition is associative). The map $x\circ\gamma$ is undergraduate differentiable by assumption; $y\circ x^{-1}$ is a map from part of $\R^d$ to part of $\R^d$ which one knows \emph{only to be continuous} (``undergraduate continuous''), as guaranteed by the definition of a topological manifold.
But the composition of a differentiable map with a merely continuous one is in general only continuous, ``a continuous sharp transition map can make an edge and destroy the differentiability''. So the strategy does not work
\ts{4}{1}{9:43}\ts{4}{1}{10:04}\ts{4}{1}{10:25}\ts{4}{1}{10:45}\ts{4}{1}{11:06}\ts{4}{1}{11:27}\ts{4}{1}{11:48}\ts{4}{1}{12:09}\ts{4}{1}{12:30}.
For a function or a map between manifolds the diagram is bigger (charts on both sides), but it is always the same problem: if one changes chart and knows only that the transition maps are continuous, one is in trouble \ts{4}{1}{12:30}\ts{4}{1}{12:50}\ts{4}{1}{13:11}.

\section{Compatible charts}\label{sec:c04-compatible}

But there is a remedy \ts{4}{1}{13:11}\ts{4}{1}{13:36}\ts{4}{1}{13:57}. In the previous section any imaginable charts were used, in other words charts from the maximal atlas $\mathcal A_{\max}$ of $(M,\mathcal O)$. Recall: an atlas is a collection of enough charts that the chart domains cover the manifold;
the maximal atlas has many more redundant pages, namely every chart one could put in (``it is a theoretical object; if I wanted to give you an atlas, I would write down a minimal or almost minimal one because it saves work''); a chart is always in the maximal atlas
\ts{4}{1}{13:57}\ts{4}{1}{14:19}\ts{4}{1}{14:43}\ts{4}{1}{15:06}\ts{4}{1}{15:28}\ts{4}{1}{15:49}. Maybe one can \emph{reduce} the maximal atlas by throwing out all the charts whose transition maps are only continuous but not differentiable, in a clever way such that what remains is still an atlas (still covers $M$) but all the
transition functions are differentiable. Then one only allows charts from this restricted atlas to define differentiability, and everything is fine \ts{4}{1}{16:09}\ts{4}{1}{16:30}\ts{4}{1}{16:50}\ts{4}{1}{17:10}.

\begin{definition}[New concepts: $\star$-compatible charts, $\star$-atlas, $\star$-manifold]\label{def:c04-compatible}
Schuller writes ``$\star$'' (a placeholder drawn as a little flower on the board [?]) for any property of maps $\R^d\to\R^d$ (e.g.\ differentiability): this is ``the philosophy that one can have different types of manifolds and different types of compatibility''.
\begin{enumerate}
  \item Two charts $(U,x)$ and $(V,y)$ of a topological manifold are called \textbf{$\star$-compatible} if either $U\cap V=\emptyset$ (``somebody you don't meet you cannot be incompatible with''), or $U\cap V\neq\emptyset$ and the chart transition maps
        \[
          y\circ x^{-1}\colon x(U\cap V)\to y(U\cap V)\quad\text{and}\quad x\circ y^{-1}\colon y(U\cap V)\to x(U\cap V)
        \]
        have the property $\star$ as maps $\R^d\to\R^d$ (careful with the domains: $x(U\cap V)$ and $y(U\cap V)$; one is the inverse of the other); property $\star$ is a notion of undergraduate analysis (e.g.\ differentiability).
  \item A \textbf{$\star$-atlas} is an atlas $\mathcal A_\star\subseteq\mathcal A_{\max}$ (a subset of the maximal atlas that still covers $M$) any two charts of which are $\star$-compatible.
  \item A \textbf{$\star$-manifold} is a triple $(M,\mathcal O,\mathcal A_\star)$: a topological manifold together with a chosen selection of mutually $\star$-compatible charts
\end{enumerate}
\ts{4}{1}{17:31}\ts{4}{1}{17:52}\ts{4}{1}{18:14}\ts{4}{1}{18:35}\ts{4}{1}{18:55}\ts{4}{1}{19:17}\ts{4}{1}{19:39}\ts{4}{1}{20:02}\ts{4}{1}{20:23}\ts{4}{1}{20:45}\ts{4}{1}{21:05}\ts{4}{1}{21:32}\ts{4}{1}{21:53}\ts{4}{1}{22:14}.
\end{definition}
That is a \emph{choice}: for a given topological manifold, say the sphere, the maximal atlas could be reduced in more than one way (keeping one set of charts or another) such that both are $\star$-compatible, and the choices need not be the same. If one wishes to model spacetime, one should know which
choice one makes, because under one choice some curves will be differentiable and under another they will not. The question of \emph{how many} choices there are is therefore very important, ``not mathematical sophistication for its own sake, a really down-to-earth question for physics: am I choosing the right structure?''
\ts{4}{1}{22:14}\ts{4}{1}{22:34}\ts{4}{1}{22:55}\ts{4}{1}{23:15}\ts{4}{1}{23:36}\ts{4}{1}{23:56}.

\paragraph{Student question: why include the empty overlap?} Why is the case of empty overlap part of the definition of compatibility? Because charts with disjoint domains can never cause an incompatibility: a curve in the world which hits both domains is judged by one chart in one part and by the other chart in the other part, never in conflict;
only where the domains intersect can the same piece of the curve be judged in two charts, and there the condition has to be imposed \ts{4}{1}{23:56}\ts{4}{1}{24:16}\ts{4}{1}{24:36}\ts{4}{1}{24:57}\ts{4}{1}{25:18}.

\subsection{What can $\star$ be}

The table of ``what can $\star$ be'' \ts{4}{1}{25:18}:
\begin{itemize}
  \item $\star=C^0$: the transition maps are continuous with respect to the standard topology on $\R^d$ (``silly, but'') \ts{4}{1}{25:38}\ts{4}{1}{25:58}. Every atlas on a topological manifold is a $C^0$-atlas, because the transition maps are continuous \ts{4}{1}{25:58}\ts{4}{1}{26:18}.
  \item $\star=C^1$: the transition maps $\R^d\to\R^d$ (strictly, from where they are defined to where they are defined) are differentiable \emph{and the result (the derivative) is continuous}; one could have a function that is differentiable without the derivative being continuous \ts{4}{1}{26:18}\ts{4}{1}{26:39}\ts{4}{1}{27:01}.
        Generalised: $C^k$, $k$ times continuously differentiable \ts{4}{1}{27:01}\ts{4}{1}{27:26}.
        Why insist on continuity of the result? One does not have to, but the nice thing about maps whose derivatives are continuous is that continuous differentiability can be judged from the existence and continuity of the partial derivatives, whereas there are maps which are not $k$ times differentiable as maps $\R^d\to\R^d$ although all the partial derivatives exist;
        ``it is nicer not to have to worry about these cases'', hence this is a more favoured class, easier to work with \ts{4}{1}{27:26}\ts{4}{1}{27:47}\ts{4}{1}{28:07}\ts{4}{1}{28:27}.
  \item $\star=C^\infty$: continuously differentiable arbitrarily many times, ``smooth'' \ts{4}{1}{28:27}\ts{4}{1}{28:48}.
  \item $\star=C^\omega$: real analytic, i.e.\ the transition map can be Taylor expanded (a multi-dimensional Taylor expansion). ``In physics sometimes people say: everything can be Taylor expanded. No, it is simply not true''; this is a much smaller class than $C^\infty$.
        A student asks whether it matters if one has an analytic function in one or in two dimensions: they are totally different functions, one cannot even compare. Smooth functions need not be analytic, in any dimension \ts{4}{1}{28:48}\ts{4}{1}{29:10}\ts{4}{1}{29:31}\ts{4}{1}{30:54}\ts{4}{1}{31:15}.
  \item ``Complex differentiable'' (for even-dimensional manifolds): the chart maps $\R^{2d}\to\R^{2d}$ pairwise satisfy the Cauchy--Riemann equations. Then one gets a \emph{complex manifold}, which has much more structure than a real manifold because of the additional equations, always asked of the transition functions. So in the future, hearing of a
        complex manifold, one knows: the Cauchy--Riemann equations apply to the transition maps \ts{4}{1}{29:31}\ts{4}{1}{29:51}\ts{4}{1}{30:12}\ts{4}{1}{31:15}\ts{4}{1}{31:35}.
\end{itemize}
Why impose more and more conditions on the transition functions? Because one may want to define a Taylor-expandable curve, and then one needs an analytic manifold in order not to destroy Taylor-expandability in one chart when passing to another: the more fancy things one wants of objects on the manifold, the more restrictive one has to be in the choice of atlas
\ts{4}{1}{30:12}\ts{4}{1}{30:33}.

So, a curve on the manifold can be called $k$ times continuously differentiable when one manages to find a $C^k$-atlas, since that preserves the property from one chart to another \ts{4}{1}{31:35}\ts{4}{1}{31:55}.

\subsection{Whitney's theorem and smooth manifolds}

To what degree do derivatives appear in physics? Is $C^2$ enough, maybe $C^3$ to be safe? There is a nice theorem by Whitney, quoted without proof, that says one need not worry \ts{4}{1}{31:55}\ts{4}{1}{32:16}\ts{4}{1}{32:43}.

\begin{theorem}[Whitney]\label{thm:c04-whitney}
Any $C^k$-atlas with $k\ge1$ of a (topological) manifold contains, as a subatlas which is still an atlas, a $C^\infty$-atlas
\ts{4}{1}{32:43}\ts{4}{1}{33:04}\ts{4}{1}{33:24}\ts{4}{1}{33:45}.
\end{theorem}
(Quoted without proof; Schuller first states it for ``$C^k$ manifolds'', a topological manifold on which there exists a $C^k$-atlas with $k\ge1$, and then reformulates it in terms of atlases ``because it is shorter and clearer''.)
Meaning: for some topological manifolds one cannot remove charts in such a way that what remains is still an atlas with continuously differentiable transition functions (they cannot be given such a structure). But for those where one can achieve at least $C^1$, one can keep removing charts until the atlas is even $C^\infty$.
The difficult step is from $C^0$ to $C^1$, ``really substantial; this might or might not work for a given topological manifold''; once one arrived at $C^1$ one trickles down to $C^\infty$ \ts{4}{1}{33:45}\ts{4}{1}{34:06}\ts{4}{1}{34:27}\ts{4}{1}{34:48}\ts{4}{1}{35:08}.

Hence, without loss of generality, one will always consider $C^\infty$-manifolds (``instead of specifying which $C^k$ I will always use''), which are called \textbf{smooth manifolds}, ``nicer and smoother than saying $C^\infty$ manifold'' \ts{4}{1}{35:08}\ts{4}{1}{35:31}.
This is unless one wishes to define Taylor expandability, complex differentiability and so on, because that does not trickle down; but one does not need it, although physicists might like to; then one has to make sure that the underlying topological manifold still allows the reduction of the atlas: ``as a physicist you cannot say, come on, let us Taylor expand a little''
\ts{4}{1}{35:31}\ts{4}{1}{35:53}\ts{4}{1}{36:13}.

\begin{definition}[New concepts: smooth manifold]\label{def:c04-smooth}
A \textbf{smooth manifold} is a triple $(M,\mathcal O,\mathcal A)$, where $(M,\mathcal O)$ is a topological manifold and $\mathcal A$ is a $C^\infty$-atlas (a subset of the maximal atlas that still covers $M$, all of whose transition maps are $C^\infty$)
\ts{4}{1}{36:35}\ts{4}{1}{36:58}.
\end{definition}
What distinguishes a smooth manifold from a topological manifold is the provision of \emph{extra structure}: not extra in the sense that something is put on top, but because the atlas has been reduced. Ripping out pages (7, 28, 510, \dots) also contains \emph{information}, ``because I could have ripped out other pages, and send the message: I love you, or I don't quite love you''
(recall the maximal atlas of \cref{def:c02-maximal}) \ts{4}{1}{36:58}\ts{4}{1}{37:19}.

\subsection{Student questions: the chalk, and curves that are only $C^2$}

\paragraph{Where does $\gamma$ come from in the experiment?} A student asks where the curve $\gamma$ comes from if one starts from physics. Schuller throws a piece of chalk. It makes a trajectory in the real world: the curve $\gamma$ made by the real object (with sufficient idealisation, with a parameter that one thinks of as time, ``later on we won't'').
To communicate to the person who did not look but was looking out of the window, one sets up a chart (``I put up little flags everywhere''): ``it passed my nose right here, at that time'', and writes it down, maybe as a nice function. But as a physicist, the trajectory of the chalk is not the chart representative $x\circ\gamma$, it is $\gamma$.
When writing down physics on paper one almost always talks about the chart representative, and may change chart and talk about another one, but what the model of the world should attach properties to is the real path of the real chalk. If one has the fantasy that ``the curve I wrote down to you was differentiable'', that is true
only because very special provisions were made of not choosing charts that do not allow one to speak of that \ts{4}{1}{37:39}\ts{4}{1}{38:06}\ts{4}{1}{38:27}\ts{4}{1}{38:48}\ts{4}{1}{39:08}\ts{4}{1}{39:29}\ts{4}{1}{39:49}.

\paragraph{What if the fifth derivative is discontinuous?} A student: what about a pulse, a hammer on a nail, so that for instance the fifth derivative of the trajectory is discontinuous? Does a smooth manifold exclude it? No: on a smooth manifold one can still define $C^2$ curves
(``if I say I want this to be $C^2$'', i.e.\ one does not require the curve to be $C^\infty$ but only $C^2$). One can do this on a smooth manifold with the $C^\infty$ transition functions, because $C^\infty$ after $C^2$ is generically only $C^2$; it does not harm that one chooses a $C^\infty$ manifold if one still wants to consider, for whatever reason,
curves that are only $C^2$ but not $C^3$. Whitney's theorem does not say that the objects one looks at must then be $C^\infty$, but that it does not harm the $C^2$ objects if one chooses the manifold $C^\infty$, and one always can. ``It is a good question, because maybe I overstated before: I said that we can always choose our objects to be smooth;
we can still insist on them not being smooth, but for the chart transition maps I still require smoothness and everything will be well defined'' \ts{4}{1}{39:49}\ts{4}{1}{40:10}\ts{4}{1}{40:30}\ts{4}{1}{40:51}\ts{4}{1}{41:12}\ts{4}{1}{41:33}\ts{4}{1}{41:54}\ts{4}{1}{42:14}.

\section{Diffeomorphisms}\label{sec:c04-diffeo}

The last section \ts{4}{1}{42:14}\ts{4}{1}{42:35}. Whenever one has maps from one set to another, and the sets have no further structure, one studies them by maps between them, and what one requires is that the map is invertible.
For ``naked'' sets $M,N$, the structure-\emph{preserving} maps are the bijections (at the set-theoretical level a set consists of elements; if each element on one side can be mapped to one on the other, everything is hit and none twice, one ``pairs girls with boys''). Two sets may or may not admit such a map;
for example between $\{1,2,3\}$ and $\{A,B\}$ there is none: one can map $1\mapsto A$, $2\mapsto B$ but then $3$ goes to $A$ or $B$ again and the map is not bijective, no inverse exists. One compares the sizes of sets by the existence of a bijection (a little overdone for finite sets, but it works for infinite ones)
\ts{4}{1}{42:35}\ts{4}{1}{42:56}\ts{4}{1}{43:17}\ts{4}{1}{43:37}\ts{4}{1}{43:58}\ts{4}{1}{44:19}\ts{4}{1}{44:39}\ts{4}{1}{45:00}.

\begin{definition}[New concepts: set-theoretically isomorphic]\label{def:c04-setiso}
Two sets $M$ and $N$ are called \textbf{(set-theoretically) isomorphic} if there exists a bijection $\phi\colon M\to N$ (such a map is then called an \textbf{isomorphism}; there are usually many different ones; ``iso'' is the same, ``morph'' is form: the same form, looking only at the set-theoretic aspects)
\ts{4}{1}{45:00}\ts{4}{1}{45:21}\ts{4}{1}{45:42}\ts{4}{1}{46:02}.
\end{definition}
Examples: $\mathbb Z$ is set-theoretically isomorphic to $\mathbb N$ (as many integers as natural numbers; one says there are double as many, apart from zero, but one can make pairs); more surprisingly the fractions are as many as $\mathbb N$ (the diagonal counting scheme); but $\mathbb N$ is not set-theoretically isomorphic to $\R$
\ts{4}{1}{46:02}\ts{4}{1}{46:22}\ts{4}{1}{46:44}.

Now add structure to the set, a topology, as in the first lecture. Is there a meaningful notion of \emph{topological} isomorphism? There is a household name: \textbf{homeomorphic}, traditional language. $M$ and $N$ are homeomorphic if they are isomorphic at the set level and in addition there exists a bijection $\phi\colon M\to N$ such that
$\phi$ \emph{and its inverse are continuous}, since continuity is the structure-respecting property here. If one exists (there may be several) the spaces are called topologically equivalent; if two spaces are homeomorphic then they are also set-theoretically isomorphic: a stronger notion, ``I look at the sets more closely''
\ts{4}{1}{46:44}\ts{4}{1}{47:04}\ts{4}{1}{47:25}\ts{4}{1}{47:47}\ts{4}{1}{48:08}\ts{4}{1}{48:29}\ts{4}{1}{48:49}.
For vector spaces the same: $V$ is (vector space) isomorphic to $W$ if there exists a bijection $\phi$ between them which is, in both directions, linear, linearity being the structure-preserving property (\cref{def:c03-linear})
\ts{4}{1}{48:49}\ts{4}{1}{49:10}\ts{4}{1}{49:31}. ``This is always the same principle.''

\begin{definition}[New concepts: diffeomorphic, $C^\infty$ map between manifolds]\label{def:c04-diffeo}
Two $C^\infty$-manifolds $(M,\mathcal O_M,\mathcal A_M)$ and $(N,\mathcal O_N,\mathcal A_N)$ are said to be \textbf{diffeomorphic} if there exists a bijection $\phi\colon M\to N$ (a purely set-theoretic thing) such that $\phi$ and $\phi^{-1}$ are both $C^\infty$-maps. Such a $\phi$ is a \textbf{diffeomorphism}.
\ts{4}{1}{49:31}\ts{4}{1}{49:52}\ts{4}{1}{50:13}\ts{4}{1}{50:34}

The notion of a $C^\infty$ map between manifolds is not yet defined, but it is clear how that goes: for $M$ ($d$-dimensional) and $N$ ($e$-dimensional) choose an open set $U\subseteq M$, and in the target an open $V\subseteq N$, and look only at that portion of the map. Choose a chart $x$ on $U$ (to $\R^d$) and a chart $y$ on $V$ (to $\R^e$);
the portion of $\phi$ is called $C^\infty$ if
\[
  y\circ\phi\circ x^{-1}\colon\R^d\supseteq x(U)\to y(V)\subseteq\R^e
\]
is undergraduate $C^\infty$, a map that uses two charts for its definition \ts{4}{1}{50:34}\ts{4}{1}{50:54}\ts{4}{1}{51:15}\ts{4}{1}{51:35}\ts{4}{1}{51:57}.
\end{definition}
As always one has to check that the definition does not depend on the charts: with alternative charts $\tilde x,\tilde y$ the representative is $\tilde y\circ\phi\circ\tilde x^{-1}=(\tilde y\circ y^{-1})\circ(y\circ\phi\circ x^{-1})\circ(x\circ\tilde x^{-1})$, and since the manifolds were taken $C^\infty$,
the two chart transition maps are undergraduate $C^\infty$ and this is preserved: the diagram gets bigger, but the principle is always the same \ts{4}{1}{51:57}\ts{4}{1}{52:18}\ts{4}{1}{52:39}.

Why bother about diffeomorphisms? If one looks at a manifold closely enough to see that it is a $C^\infty$ manifold, can one distinguish two $C^\infty$ manifolds using only that structure? Example: take the unit sphere $\{(m,n,p)\in\R^3\mid m^2+n^2+p^2=1\}$, embedded, with the topology pulled back from the standard topology of $\R^3$; one gets a topological manifold,
and one can reduce the maximal atlas to a $C^\infty$ one (``I promise you you can''). Do the same with $5m^2+7n^2+18p^2=1$, an ellipsoid. Is there a diffeomorphism between them? Yes. Between the sphere and a sphere ``where I make a little edge in it'', folded a little? No: there is a homeomorphism (as topological manifolds they are the same),
but as smooth manifolds they are not the same any more. ``It always depends at the level you look at it.'' The surface of a potato without folds is diffeomorphic to the surface of the Earth: at the level of smooth manifolds things do not have a \emph{shape} yet, apart from not having folds or edges; the shape comes only later
\ts{4}{1}{52:39}\ts{4}{1}{52:59}\ts{4}{1}{53:19}\ts{4}{1}{53:40}\ts{4}{1}{54:00}\ts{4}{1}{54:22}\ts{4}{1}{54:42}\ts{4}{1}{55:03}.

\subsection{How many smooth structures are there?}

The last thing: a theorem on the number of $C^\infty$-manifolds one can make out of a given $C^0$-manifold (if any: for some it might not work), \emph{up to diffeomorphism}: if two manifolds differ only by a diffeomorphism one counts them as one, since from the differentiable-manifold perspective they are the same
\ts{4}{1}{55:03}\ts{4}{1}{55:28}\ts{4}{1}{55:49}\ts{4}{1}{56:09}.

\begin{theorem}[Moise, Radon: number of smooth structures, quoted]\label{thm:c04-count}
Let $(M,\mathcal O)$ be a topological manifold of dimension $d$ that admits a smooth structure. Up to diffeomorphism:
\begin{enumerate}
  \item for $d=1,2,3$ there is exactly one smooth structure (``one, one, one: no ambiguity'');
  \item for $d=4$: open (see below);
  \item for $d=5,6,7,\dots$ there are finitely many (``five, six, seven, finitely many''); as announced, this holds in each such dimension, though the number may depend on the manifold
\end{enumerate}
\ts{4}{1}{56:09}\ts{4}{1}{56:30}\ts{4}{1}{56:52}\ts{4}{1}{57:12}.
\end{theorem}
(Schuller calls it ``the Moise--Radon theorem'': a mishearing or a conflation of Moise and Rad\'o [?]; see the history box.) This is not easy to prove. That there are only finitely many in dimensions $5,6,7,\dots$ was shown by a field of topology called \emph{surgery theory}, where one has to learn how to glue together manifolds and so on: ``extra techniques had to be invented
in topology to prove this, so don't worry if you don't feel it; it is a deep result'' \ts{4}{1}{56:52}\ts{4}{1}{57:12}.

Now four dimensions, ``which of course, as spacetime physicists, we are particularly interested in''. There are some partial results: for some compact manifolds (compactness is a topological notion; the surface of a sphere is compact but $\R^2$ is not) with, say, Betti number $18$ (also a topological notion)
there are only finitely many choices or something like this, ``some very very special exceptions''; but generically, and particularly for \emph{non-compact} manifolds (``which is a real option for spacetime''), there exist \emph{uncountably infinitely many}, ``not only infinitely many 1, 2, 3, 4, 5, but so many like the real numbers that you can't even count them'':
there exist uncountably many different $C^\infty$-manifolds one can make out of a given topological manifold \ts{4}{1}{57:12}\ts{4}{1}{57:33}\ts{4}{1}{57:53}\ts{4}{1}{58:14}\ts{4}{1}{58:34}. ``That is just an existence proof. So that is a real problem: which one do you choose for physics?''

\paragraph{Student question: is the finite number the same for all?} No. ``How finite? Two, three or ten? That does not matter. I don't know whether the proofs say anything about a definite number; probably not, because these could look very funny: some five-dimensional ones might have 20 and other five-dimensional manifolds might
have 28 different structures; I would be very surprised if in dimension 28 the number was 58 for all, it could be, I don't know; but finite still sounds fine.'' For a physicist one can say ``let experiment decide''; but then one would be a very busy experimental physicist: ``you see, theory does something here for you; you cannot always say let experiment decide''
\ts{4}{1}{58:34}\ts{4}{1}{58:55}\ts{4}{1}{59:15}\ts{4}{1}{59:36}\ts{4}{1}{59:56}\ts{4}{1}{1:00:17}.

\begin{history}[Smooth structures by dimension]
Schuller attributes the dimension $\le3$ result to ``Moise and Radon''. \emph{Checked.} The topological dimension $2$ case goes back to Rad\'o (1925: every surface can be triangulated) and the dimension $3$ case to Moise (1952, triangulation and the Hauptvermutung for $3$-manifolds); together with Munkres and Whitehead
they give the uniqueness of the smooth structure up to diffeomorphism in dimensions $\le3$. In dimension $4$, $\R^4$ itself carries uncountably many pairwise non-diffeomorphic smooth structures (``exotic $\R^4$'', Taubes 1987, building on Donaldson and Freedman), which is the result alluded to for non-compact manifolds.
Whitney's theorem is from his 1936 paper on differentiable manifolds. (Not said in class.)
\end{history}
\begin{coursenote}[Finite in dimension $\ge5$]
The finiteness for $d\ge5$ is, as far as it is known to the note-taker, a statement for \emph{compact} manifolds (surgery theory); Schuller states it for all topological manifolds of that dimension admitting a smooth structure. (Not said in class; to be checked [?].)
\end{coursenote}

% ==================================================================
\section{Tutorial 4}\label{sec:t04-1}

Tutorial 4, 0:35:23, \seg{4}{2}{video}, by a tutor, not Schuller. The sheet (Sheet 4) is paraphrased, not copied; see \url{https://tales.mbivert.com/on-heraeus-winter-school-tutorials/} (and \url{https://tales.mbivert.com/heraeus/Sheet_04.pdf}).
The first exercise is again fifteen true/false statements, in three groups of five \ts{4}{2}{0:00}\ts{4}{2}{0:21}.

\subsection*{Exercise 1: true or false}

\paragraph{(A) Differentiability classes.} $C^k$ is the set of functions that are $k$ times continuously differentiable.
\begin{itemize}
  \item $f\colon\R\to\R$, $f(x)=x^2$ lies in $C^3$: \textbf{true}; the derivatives are $2x$, $2$, $0$, \dots, all continuous \ts{4}{2}{0:21}\ts{4}{2}{0:41}\ts{4}{2}{1:02}.
  \item $f(x)=|x|$ lies in $C^2$: \textbf{false}; there is a kink at $0$, so $f$ is not even differentiable there \ts{4}{2}{1:02}\ts{4}{2}{1:22}.
  \item $f(x)=|x|^3$ lies in $C^3$: \textbf{false}; $f'=3x^2$ for $x>0$, $-3x^2$ for $x<0$; $f''=6|x|$, which has a kink at $0$, so the third derivative does not exist there (the function is $C^2$ but not $C^3$) \ts{4}{2}{1:44}\ts{4}{2}{2:06}\ts{4}{2}{2:26}.
  \item $f(x)=\exp|x|$ lies in $C^\infty$: \textbf{false}; again a kink at $0$ \ts{4}{2}{2:26}\ts{4}{2}{2:47}.
  \item $f(x)=\ln x$ on $\R^+$ lies in $C^\infty$: \textbf{true} \ts{4}{2}{3:08}.
\end{itemize}

\paragraph{(B) Differentiable manifolds.}
\begin{itemize}
  \item Every smooth manifold has a maximal atlas: \textbf{true}; a smooth manifold has an atlas that covers it with conditions on the transition functions, and any atlas can be extended to a maximal one \ts{4}{2}{3:32}\ts{4}{2}{3:53}.
  \item Every $C^2$ manifold is also a $C^1$ manifold: \textbf{true}, every $C^2$ function is $C^1$ \ts{4}{2}{3:53}\ts{4}{2}{4:15}.
  \item Every differentiable manifold is also a topological manifold: \textbf{true}; ``differentiable'' means $C^k$ with $k\ge1$, ``topological'' means $C^0$ \ts{4}{2}{4:15}\ts{4}{2}{4:35}.
  \item In a $C^\infty$ manifold the chart maps are differentiable: \textbf{true}. A function $f\colon U\to\R^d$ on a chart domain is differentiable if its chart representative $f\circ\tilde x^{-1}$ is, for any chart $\tilde x$ of the smooth atlas; taking $f=x$ and $\tilde x=x$ one gets $x\circ x^{-1}=\mathrm{id}_{\R^d}$, clearly differentiable
        \ts{4}{2}{4:35}\ts{4}{2}{4:56}\ts{4}{2}{5:17}\ts{4}{2}{5:38}\ts{4}{2}{5:59}\ts{4}{2}{6:19}.
  \item The definition of a $C^k$ curve on a manifold requires a $C^m$ manifold with $m\ge k$: \textbf{true}; the atlas has to be at least $C^k$ to judge $C^k$ differentiability \ts{4}{2}{6:19}\ts{4}{2}{6:40}.
\end{itemize}

\paragraph{(C) Further statements.}
\begin{itemize}
  \item A differentiable manifold is always a topological space: \textbf{true}; adding an atlas keeps the topological space, which gets more structure by becoming a manifold \ts{4}{2}{7:01}\ts{4}{2}{7:23}.
  \item A differentiable manifold has enough structure to constitute a vector space: \textbf{false} \ts{4}{2}{7:23}.
  \item A differentiable manifold is a generalisation of a set: \textbf{false}; it is a specification of a set (a set with additional structure) \ts{4}{2}{7:23}\ts{4}{2}{7:44}.
  \item A differentiable manifold requires the discrete topology: \textbf{false}; it can carry any topology it likes (compatible with the charts) \ts{4}{2}{7:44}.
  \item A differentiable manifold features at least one set of charts that covers the whole manifold: \textbf{true}, the idea of an atlas \ts{4}{2}{7:44}\ts{4}{2}{8:05}.
\end{itemize}

\subsection*{Exercise 2: an atlas on $\R$ that is not smooth}

Consider $\R$ with the standard topology and the atlas $\mathcal A=\{(\R,x),(\R,y)\}$ with the identity $x=\mathrm{id}$ and $y(a)=a^3$; each chart covers all of $\R$ \ts{4}{2}{8:05}\ts{4}{2}{8:26}\ts{4}{2}{8:47}.
\begin{enumerate}
  \item The transition map $y\circ x^{-1}\colon a\mapsto a^3$ is $C^\infty$ \ts{4}{2}{9:08}\ts{4}{2}{9:28}\ts{4}{2}{9:49}\ts{4}{2}{10:10}.
  \item The other direction $x\circ y^{-1}\colon a\mapsto\sqrt[3]a$ is continuous but not differentiable at $0$ (the tangent of the graph is parallel to the vertical axis, the derivative has a pole): not even $C^1$, only $C^0$ \ts{4}{2}{10:10}\ts{4}{2}{10:30}\ts{4}{2}{10:52}\ts{4}{2}{11:16}\ts{4}{2}{11:38}.
        So with this atlas $\R$ is \emph{not} a differentiable manifold \ts{4}{2}{11:38}.
  \item Repair: remove one of the two charts (it does not matter which). The remaining single chart still covers $\R$, there are no transition maps, hence the atlas is trivially smooth: $\R$ with the standard topology and this atlas is a smooth manifold. This is what was said in the lecture:
        restricting an atlas of a topological manifold can make it differentiable \ts{4}{2}{12:00}\ts{4}{2}{12:21}\ts{4}{2}{12:43}\ts{4}{2}{13:04}.
\end{enumerate}

\subsection*{Exercise 3: an atlas of squares on $\R^2$}

On $\R^2$ with the ``soft square'' topology [?, the topology generated by open squares, which is the standard topology of $\R^2$] consider the sets $U_n:=\{(a,b)\in\R^2\mid|a|<n,\ |b|<n\}$, $n\in\mathbb N^+$, with the maps
\[
  x_n\colon U_n\to\R^2,\qquad x_n(a,b)=\Bigl(\frac{a+b}{2n},\,\frac{a-b}{2n}\Bigr)
\]
\ts{4}{2}{13:04}\ts{4}{2}{13:24}\ts{4}{2}{13:44}\ts{4}{2}{14:06}\ts{4}{2}{14:26}.
\begin{enumerate}
  \item \emph{The $(U_n,x_n)$ are charts.} Each $U_n$ is open (squares generate the topology). The map $x_n$ is continuous, and invertible with inverse $x_n^{-1}\colon x_n(U_n)\to U_n$, $(a,b)\mapsto(n(a+b),\,n(a-b))$, also continuous \ts{4}{2}{14:26}\ts{4}{2}{14:47}\ts{4}{2}{15:09}\ts{4}{2}{15:29}\ts{4}{2}{15:50}\ts{4}{2}{16:11}.
  \item \emph{A $C^k$ atlas, and which $k$.} The transition map from the chart $x_n$ to the chart $x_m$ is on $x_n(U_n\cap U_m)$
        \[
          x_m\circ x_n^{-1}(a,b)=x_m\bigl(n(a+b),n(a-b)\bigr)=\frac nm\,(a,b)
        \]
        (the tutor first gets $m/n$, then corrects to $n/m$). Since $m,n$ are arbitrary this is the most general transition map, and it is a multiplication by a constant, hence $C^\infty$: the atlas is a smooth atlas and $\R^2$ with it a smooth manifold, $k=\infty$
        \ts{4}{2}{16:33}\ts{4}{2}{16:55}\ts{4}{2}{17:15}\ts{4}{2}{17:37}\ts{4}{2}{18:01}\ts{4}{2}{18:22}\ts{4}{2}{18:45}\ts{4}{2}{19:07}\ts{4}{2}{19:27}\ts{4}{2}{19:49}.
  \item \emph{Another chart in the maximal extension.} The maximal extension of $\mathcal A$ contains every chart that is $C^\infty$-compatible with all those of $\mathcal A$. Take $(U_5,\tilde x)$ with $\tilde x=\mathrm{id}$ on the domain $U_5$ (``it does not really matter that I picked $U_5$'').
        The transition maps are
        \[
          \tilde x\circ x_n^{-1}(a,b)=\bigl(n(a+b),\,n(a-b)\bigr),\qquad x_n\circ\tilde x^{-1}(a,b)=\Bigl(\frac{a+b}{2n},\frac{a-b}{2n}\Bigr),
        \]
        both smooth (linear), hence the chart lies in the maximal extension of $\mathcal A$ (the tutor first writes the two lines the wrong way round and corrects himself)
        \ts{4}{2}{20:13}\ts{4}{2}{20:34}\ts{4}{2}{20:54}\ts{4}{2}{21:15}\ts{4}{2}{21:38}\ts{4}{2}{21:59}\ts{4}{2}{22:21}\ts{4}{2}{22:41}\ts{4}{2}{23:02}.
\end{enumerate}

\subsection*{Exercise 4: partial derivatives as maps}

A short exercise, but important because it introduces a good notation. For $f\colon\R^d\to\R$ the symbol $\partial_if$ denotes the partial derivative with respect to the $i$-th \emph{entry}; the important point is that $\partial_if$ is again a map in its own right, $\R^d\to\R$
\ts{4}{2}{23:23}\ts{4}{2}{23:44}\ts{4}{2}{24:04}. For a given $f\colon\R^3\to\R$ [the sheet's function is not read out; from the results it is $f(\alpha,\beta,\delta)=\alpha^3\beta^2+\beta^2\delta+\delta$, reconstructed [?]] one computes, e.g.:
\begin{itemize}
  \item $\partial_2f=2\alpha^3\beta+2\beta\delta+0$, as a map evaluated at $(x,y,z)$: $(\partial_2f)(x,y,z)=2x^3y+2yz$;
  \item $\partial_1f=3\alpha^2\beta^2$, evaluated at $(x,y,z)$: $3x^2y^2$; ``really one first takes the derivative and then evaluates at some place'';
  \item second derivatives: $\partial_1\partial_2f=6\alpha^2\beta$, evaluated at $(a,b,c)$: $6a^2b$;
  \item $\partial_3\partial_3f$: the first derivative is $\beta^2+1$ and the second is zero, ``not because we plugged in the numbers first and then derived some constants, which are zero, but because the second derivative itself is the zero function''
\end{itemize}
\ts{4}{2}{24:04}\ts{4}{2}{24:46}\ts{4}{2}{25:28}\ts{4}{2}{26:10}\ts{4}{2}{26:53}.

\subsection*{Exercise 5: functions and curves in a chart}

``How to deal with functions and curves in a chart''; visualising these maps with a diagram always helps \ts{4}{2}{26:53}\ts{4}{2}{27:13}. Let $(M,\mathcal O,\mathcal A)$ be a $d$-dimensional smooth manifold, $(U,x)\in\mathcal A$ a chart, $\gamma\colon\R\to U$ a smooth curve and $f\colon U\to\R$ a smooth function
\ts{4}{2}{27:13}\ts{4}{2}{27:34}.
\begin{enumerate}
  \item \emph{The commutative diagram.} Draw the manifold with the subset $U$, the curve $\gamma\colon\R\to U$, the function $f\colon U\to\R$ and the chart $x\colon U\to x(U)\subseteq\R^d$. What is missing are the chart representatives, i.e.\ everything pushed down into the chart: the representative of the curve is $x\circ\gamma\colon\R\to\R^d$, and the representative of the function
        is $f\circ x^{-1}\colon x(U)\to\R$ (one goes up with $x^{-1}$ and then applies $f$) \ts{4}{2}{27:34}\ts{4}{2}{27:55}\ts{4}{2}{28:16}\ts{4}{2}{28:36}\ts{4}{2}{28:57}\ts{4}{2}{29:18}\ts{4}{2}{29:39}.
  \item \emph{The derivative $\frac{\mathrm d}{\mathrm d\lambda}(f\circ\gamma)$.} Does it make sense? $f\circ\gamma$ is a map $\R\to\R$, so differentiation makes sense. To compute it one cannot use the chain rule directly, because $f$ and $\gamma$ have domain and target on the manifold, where there is no notion of differentiation. Therefore go through the chart representatives:
        \[
          f\circ\gamma=(f\circ x^{-1})\circ(x\circ\gamma)=:F\circ\Gamma
        \]
        (inserting $x^{-1}\circ x=\mathrm{id}$, composition being associative), with $F:=f\circ x^{-1}$ and $\Gamma:=x\circ\gamma$. The chain rule then reads
        \[
          (f\circ\gamma)'(\lambda)=\partial_iF\bigl(\Gamma(\lambda)\bigr)\cdot(\Gamma^i)'(\lambda)
        \]
        (summation over $i$), the outer function evaluated at the inner function times the derivative of the $i$-th component of the inner function.
        Concretely, for $F(x,y)=x^2+y^2$ and $\Gamma(\lambda)=(\cos\lambda,\sin\lambda)$ one gets $2\cos\lambda\cdot(-\sin\lambda)+2\sin\lambda\cdot\cos\lambda=0$.
        This is clear geometrically: $\Gamma$ describes a circle which is run around with increasing $\lambda$, and $F=x^2+y^2$ is constant ($=1$) on the circle, so the derivative along the curve is zero
        \ts{4}{2}{30:00}\ts{4}{2}{30:20}\ts{4}{2}{30:42}\ts{4}{2}{31:05}\ts{4}{2}{31:26}\ts{4}{2}{31:47}\ts{4}{2}{32:07}\ts{4}{2}{32:28}\ts{4}{2}{32:49}\ts{4}{2}{33:09}\ts{4}{2}{33:30}\ts{4}{2}{33:51}\ts{4}{2}{34:13}\ts{4}{2}{34:34}\ts{4}{2}{34:55}.
\end{enumerate}

% ==================================================================
\section*{Key concepts}
\begin{keyconcepts}
\textbf{Retrieval questions.}
\begin{enumerate}
  \item Why is a topological manifold not enough to define differentiability of a curve? Where does the strategy via chart representatives fail?
  \item Define $\star$-compatible charts, a $\star$-atlas and a $\star$-manifold. Why are charts with disjoint domains always compatible?
  \item List the choices of $\star$ ($C^0$, $C^k$, $C^\infty$, $C^\omega$, complex) and what each requires of the transition maps. Why is a complex manifold ``more structure''?
  \item State Whitney's theorem. Which step is the substantial one, and why can one always work with smooth manifolds?
  \item Why can one still consider $C^2$ curves on a smooth manifold?
  \item Define set-theoretic isomorphism, homeomorphism, vector space isomorphism and diffeomorphism. What is the structure each one preserves?
  \item How is a map between two smooth manifolds called $C^\infty$, and why is this independent of the charts?
  \item How many smooth structures does a topological manifold carry up to diffeomorphism in dimension $1,2,3$, $5,6,7,\dots$ and $4$? Why is this a problem for physics?
\end{enumerate}
\textbf{Concepts.}
\begin{itemize}
  \item Failure of the lift: $y\circ\gamma=(y\circ x^{-1})\circ(x\circ\gamma)$ with $y\circ x^{-1}$ only continuous (\cref{sec:c04-strategy}).
  \item $\star$-compatible charts, $\star$-atlas, $\star$-manifold (\cref{def:c04-compatible}); the classes $C^0,C^k,C^\infty,C^\omega$, complex.
  \item Whitney: a $C^k$-atlas ($k\ge1$) contains a $C^\infty$-subatlas (\cref{thm:c04-whitney}); smooth manifold $(M,\mathcal O,\mathcal A)$ (\cref{def:c04-smooth}).
  \item Isomorphisms at each level of structure; diffeomorphism and $C^\infty$ maps via chart representatives $y\circ\phi\circ x^{-1}$ (\cref{def:c04-diffeo}).
  \item Number of smooth structures up to diffeomorphism: $1$ for $d\le3$, finite for $d\ge5$, uncountably many in general for $d=4$ (\cref{thm:c04-count}).
  \item Chart representatives of curves and functions and the chain rule in a chart (\cref{sec:t04-1}).
\end{itemize}
\end{keyconcepts}

\flushside
\end{paracol}
